% a01_H0rel part 6: the upper bound, the assembly, the clause (R_0)
\subsection{The upper bound}\label{app:a01:upper}

We follow \cite[Sects.~2--3]{Ba2} and \cite[Sect.~5.2]{K2}. Ba{\l}aban's six large-field families (the tests on $|B-QA|$, $|\partial A|$, $|A|$, $|\psi-Q(A)\phi|$, $|D_A\phi|$, $|\phi|$ of \cite[(2.2)]{Ba2}) are used at every scale $i\le K_f$, with the face family of \cref{def:a01:facefam} as a seventh family at every scale, $i=0$ included; at the stopping scale $j$ the small-field term carries $\chi_j\lambda_j$. The large-field decomposition, its regions $\Lambda_0^{(i)}\supset\Lambda_1^{(i)}\supset\dots\supset\Lambda_7^{(i)}$ (layers of thickness $r(\ell_i)$), the nesting rules and the inductive representation \cite[(2.43)--(2.52)]{Ba2} depend only on the set of event blocks and the number of event types, so they hold verbatim with the face family, which is a pointwise condition on one block variable at one block of the index set $\Lone^{(i)}\cap\Lambda_7^{(i-1)\prime}$. Four steps of Ba{\l}aban's printed proofs do not hold as printed at block factor $5$ and are used in a replaced form: the cube step of \cite[p.~433]{Ba3}, the cut-off $\zeta^{(k)}$ of \cite[(2.44)]{Ba2}, the boxes in the proof of \cite[Lemma~2.4]{Ba2}, and the cube partition \cite[(3.43)--(3.45)]{Ba2}, the last replaced by a product partition into boxes of large blocks with sides in $[r+s_b,2r+3s_b]$, which has every property used in \cite[Sect.~3C]{Ba2}; the replaced objects are defined from $(L,M,R,\ell_i,\bar n)$ alone and do not depend on the matter boundary condition, and the replacements change none of the printed statements. Their block-factor-$5$ form, with proofs, follows.
% a01_H0rel part 6x1: the four replaced steps of Balaban's proofs at block factor 5 (input from a01_H0rel_p6.tex)
\paragraph{The four replaced steps at block factor $5$.}
Scale-$i$ units (the block lattice of scale $i$ has spacing $1$); $u_i\defeq1+\ln\ell_i^{-1}$, so $u_{i-1}=u_i+\ln L$, $r_i\defeq r(\ell_i)=Ru_i^{\varrho}$, $p_i=p(\ell_i)=b_0u_i^{p}$ with $\varrho=2$, $p=6$; $s_b=LM$ is the large-block side; $\rho_i\defeq(r(\ell_{i-1})-s_b)/L$ is one previous-scale corridor in scale-$i$ units (\cite[p.~578]{Ba2}); $X'$ denotes the blocks of $X$ read as points of the next lattice \cite[(1.16)]{Ba1}; $\Omega\defeq B^i(\Lambda_2^{(i-1)\prime})$ is the propagator domain and $\Omega_1\defeq B^i(\Lambda_7^{(i-1)\prime})$ the vertex set of \cite[(1.1)--(1.5)]{Ba3}; $\delta_0$, $c_{\mathrm{Ba}}$ and the depth $R_{\mathrm{Ba}}$ are the constants of \cite[Props.~2.1--2.3]{Ba1}. Everything below is a statement about Ba{\l}aban's regions, cut-offs, boxes and partitions, which are built from $(L,M,R,\ell_i,\bar n)$ and the set of event blocks alone; nothing depends on the matter boundary condition, so each statement holds on the cube exactly as on the torus.

\begin{lemma}[the supply of depth, and the criterion]\label{lem:a01:X0}
(a) $\dist(\Lambda^{(i-1)}_{m+s},(\Lambda^{(i-1)}_m)^c)>s\,r(\ell_{i-1})-(s-1)$ for $s\ge1$ (scale-$(i-1)$ units), and in scale-$i$ units $\dist(B^i(\Lambda^{(i-1)\prime}_{m+n}),B^i(\Lambda^{(i-1)\prime}_m)^c)\ge n\,r(\ell_{i-1})/L-2$ for $n\le L+1$. The hypothesis of \cite[Prop.~1]{Ba3} asks depth $>2r_i$ of $\supp\tilde A$ in $\Omega$; at its use in \cite{Ba2}, where the domain is $B^i(\Lambda_0^{(i)})$, the regions supply $\dist(\supp\tilde A,\Omega^c)>2r_i-2-M_\theta-2\eta$ (the slice of $\theta_i$ and the fine blocks), which is all that the proof consumes. (b) Let $\alpha,\beta>0$, $c\ge0$. If $\alpha L\le\beta$, then $\beta\rho_i-\alpha r_i\ge\alpha R\varrho\ln L\,u_i^{\varrho-1}-\beta s_b/L$ for every $i$, so $\beta\rho_i\ge\alpha r_i+c$ at every scale once $R\ge(c+\beta s_b/L)/(\alpha\varrho\ln L)$. If $\alpha L>\beta$, then $\beta\rho_i<\alpha r_i$ for every $\ell_i<\exp\big(1-\ln L/((\alpha L/\beta)^{1/\varrho}-1)\big)$, whatever $R$ and $c$ are. If an object of size $\theta r_{i+1}$ at scale $i+1$ must sit at depth $\beta r_i$, the containment holds at every scale iff $L\theta\le\beta$, and for $L\theta>\beta$ it fails for every $\ell_i<\exp\big(1-\ln L/(1-(\beta/(L\theta))^{1/\varrho})\big)$ (now $r_{i+1}/r_i=(1-\ln L/u_i)^\varrho$ enters). (c) Where only exponential smallness across scales is used, $e^{-\delta\beta\rho_i}\le Ce^{-(\delta\beta/L)r_i}$, which is $O(\ell_i^{\kappa})$ for every $\kappa$ since $\varrho>1$.
\end{lemma}
\begin{proof}
(a) By \cite[(2.8)]{Ba2} $(\Lambda_{m+1})^c$ contains every large block within $r$ of $(\Lambda_m)^c$, so a point of $\Lambda_{m+1}$ is at distance $>r$ from $(\Lambda_m)^c$; along a lattice geodesic from $x\in\Lambda_{m+s}$ to $(\Lambda_m)^c$ each further layer adds $>r-1$. For the primed sets, points of $B^i(y)$, $B^i(y')$ with $y\in\Lambda_{m+n}$, $y'\in\Lambda_m^c$ differ coordinatewise by $(y-y')+\tau$, $|\tau|<L$, so $|x-x'|_\infty>n\,r(\ell_{i-1})-(n-1)-L$, and dividing by $L$ gives the claim for $n\le L+1$. (b) If $\beta/L\ge\alpha$ then $\beta\rho_i-\alpha r_i\ge\alpha(r_{i-1}-r_i)-\beta s_b/L$ and $r_{i-1}-r_i=R[(u_i+\ln L)^\varrho-u_i^\varrho]\ge R\varrho\ln L\,u_i^{\varrho-1}$. If $\alpha L>\beta$, $\beta\rho_i\le(\beta/L)r_{i-1}$ and $r_{i-1}/r_i=(1+\ln L/u_i)^\varrho<\alpha L/\beta$ exactly when $u_i>\ln L/((\alpha L/\beta)^{1/\varrho}-1)$. (c) $e^{-\delta\beta\rho_i}\le e^{\delta\beta s_b/L}e^{-(\delta\beta/L)r_{i-1}}$ and $e^{-cR(1+\ln\ell^{-1})^\varrho}\le\ell^\kappa$ for small $\ell$ because $\varrho>1$.
\end{proof}

Every place where $L$ enters \cite{Ba1,Ba2,Ba3} and \cite[Thm.~3.1]{K1} is a constant, a rate, a convergent sum, a condition on $(b_0,R,M,\eps_0^{\mathrm{Ba}})$, a divisibility, or a \emph{cross-scale containment} $\beta\rho_i\ge\alpha r_i+c$, and by \cref{lem:a01:X0}(b) a containment holds at every scale iff $\alpha L\le\beta$. At $L=5$ exactly four printed steps fail this test (\cref{thm:a06:BALA}): the cube step of \cite[p.~433]{Ba3} ($\alpha=\frac32$, $\beta=5$), the regularity of the vector background \cite[Lemma~2.3]{Ba2} ($\alpha=1$, $\beta=3$), the boxes of \cite[Lemma~2.4]{Ba2} ($\alpha=4$, $\beta=5$; this one fails at every $L\ge2$), and the split \cite[(2.89)]{Ba2} ($L\theta=L$, $\beta=2$); the partition \cite[(3.43)--(3.45)]{Ba2} is a divisibility entry which fails on some lattices at $L=5$. These are gaps of the printed proofs, not counterexamples to the theorems, whose statements never mention the objects concerned.

\subparagraph{(a) The cube step.} On \cite[p.~433]{Ba3} the localisation cubes of a renormalised class, which lie in a cube of side $\ell_{\bar n}$ depending on $\bar n$ only \cite[p.~432]{Ba3}, are enclosed in concentric cubes $\square_1\subset\square$ of sides $r_i$, $3r_i$, sums of large blocks; $G_i(\Omega,\tilde B)$ is written $G_i(\square,\tilde B)+\delta G_i(\Omega,\square,\tilde B)$ with the difference treated as an external field, the regular field $\tilde B$ is replaced on $\square$ by its constant part $B_0$, $B_0$ is gauged away, and $G_i(\square,0)$ is replaced by $G_i(0)$. The step presupposes $\square\subset\Omega$ (resp.\ $\square\subset\Omega_2$): the Neumann operator $G_i(\square,\tilde B)$ is built from the bonds of $\square$, where $\tilde B$ lives only if $\square\subset\Omega$, and \cite[(2.26)]{Ba1} is stated for nested domains. Every vertex of every class lies within sup-distance $1$ of $\Omega_1\cup\supp\tilde A$ (the vertices \cite[(1.6)--(1.7)]{Ba3} are summed over $\Omega_1$; those of \cite[(1.8)--(1.11)]{Ba3} carry $g_i$, $\supp g_i\subset\Omega_1$, or $\tilde A$; the pair vertices \cite[(1.12)--(1.15)]{Ba3} carry $A'$ or $\tilde A$ along a contour inside one block). For an old vertex $x_v$ of the class the containment reads
\begin{equation}\label{eq:a01:N}
\dist(x_v,\Omega^c)>\tfrac32r_\square+2s_b+\ell_{\bar n}\ \Longrightarrow\ \square\subset\Omega\ \text{and}\ \dist(\square_1,\square^c)\ge r_\square
\end{equation}
(cubes rounded outward to large blocks; triangle inequality). With the printed $r_\square=r_i$ and a vertex anchored at $\Omega_1$, whose depth is $\ge5r(\ell_{i-1})/L-2$ by \cref{lem:a01:X0}(a), this is a containment with $\alpha=\frac32$, $\beta=5$: it holds at every scale iff $L\le3$, and at $L=5$ it fails for all small $\ell_i$; the failing vertices ($\lambda$- and $\dm$-vertices summed over $\Omega_1$, $g_i$-localised photon legs near $\partial\Omega_1$) carry no $\tilde A$-leg, so no hypothesis on $\supp\tilde A$ helps.

\begin{proposition}[the cube step with a constant cube]\label{prop:a01:P12}
Let $\ell^\sharp\defeq3\ell_{\bar n}+4N_B$, where $N_B=n_1+n_2-1$ is the number of insertions per line in the expansion of \cite[p.~433]{Ba3} to order $\bar n$ (a function of $\bar n$ only), let $r^\flat$ be $\max\{3\ell^\sharp+R_{\mathrm{Ba}}+12,\,3R_{\mathrm{Ba}}+15,\,107\}$ rounded up to a multiple of $s_b$, and $D^\flat\defeq\frac32r^\flat+2s_b+\ell^\sharp$. Run the step with concentric cubes $\square_1,\square$ of sides $r^\flat$, $3r^\flat$ centred at the component cube. If every old vertex of the class has depth $>D^\flat$ in $\Omega$ (resp.\ $\Omega_2$), the proof of \cite[pp.~433--445]{Ba3} goes through and gives \cite[(1.33)]{Ba3} with a constant of the printed form, depending on $(\alpha_0,\bar n,d,a,M)$ only. If $R\ge R^\flat\defeq D^\flat+3$ and $R\ge M_\theta+2$, then at $L=5$ every class occurring in the use of \cite[Prop.~1]{Ba3} in \cite{Ba1,Ba2} has its old vertices at depth $>D^\flat$, so \cite[Prop.~1]{Ba3} holds at $L=5$ in its printed form.
\end{proposition}
\begin{proof}
The side $r_\square$ enters \cite[pp.~433--434]{Ba3} in five places. (1) \emph{Containment:} \eqref{eq:a01:N} with $r_\square=r^\flat$. (2) \emph{Old-vertex classes lie in $\square_1$:} the expansion in $\tilde B'$ uses \cite[(3.14), (3.44), (3.45)]{Ba1} to order $\bar n$; only lines between localisations whose boxes meet stay internal \cite[p.~432]{Ba3}, so along an expanded line a new vertex advances at most $2$, and the new vertices lie within $2N_B$ of the class-component cube; the boxes of all graphs of a connected component lie in a grid cube of side $3\ell_{\bar n}$ concentric with one member's cube (``localizations \dots\ in agreement'', \cite[p.~432]{Ba3}); so the class lies in a cube of side $\ell^\sharp$, and $3\ell^\sharp\le r^\flat$. (3) \emph{The two $\delta G$ external fields are $O(1)$:} $\delta G_i(\Omega,\square,\tilde B)$ is evaluated at points of $\square_1$, at distance $\ge r^\flat\ge R_{\mathrm{Ba}}+12$ from $\square^c$, with sources in $\square_1$, so \cite[(2.26)]{Ba1} for the pair $(\square,\Omega)$ gives $\norm{h\,\delta G_i\,h'}_{1,\alpha}\le c_{\mathrm{Ba}}e^{-\delta_0r^\flat}$; for $\delta G_i(\square,\eta\Z^d,0)$ take a torus $T'\supset\square$ with side a multiple of $2s_b$ and at least $\max\{3s_b,3r^\flat+2s_b+2\}$: at zero background $G_i(T',0)$ is the periodisation of $G_i(\eta\Z^d,0)$, so $\delta G_i(\square,\eta\Z^d,0)=\delta G_i(\square,T',0)+\sum_{n\ne0}G_i(\eta\Z^d,0)(\cdot,\cdot+n\ell')$, the first term obeying \cite[(2.26)]{Ba1} on $T'$ with constants independent of $\ell'$ (the constants of \cite{Ba4} depend on $d$, $a$, $M$, $L$ only) and the second vanishing as $\ell'\to\infty$ by \cite[(2.25)]{Ba1}. Both are external fields of dimension $0$ \cite[(2.1), p.~422]{Ba3}, so an $O(1)$ bound is all that is used. (4) \emph{The constant field:} for $x\in\square$ and $c$ the centre, $|\tilde B(x)-\tilde B(c)|\le d(\frac32r^\flat+s_b)\sup_\square|\partial^\eta\tilde B|=O(r^\flat p_i)=O(p_i)$ with $B_0\defeq\tilde B(c)$, in place of the printed $O(r_ip_i)\le O(p_i^2)$; every use of this bound is an upper bound monotone in the size of $\tilde B'$ (vertices weighted by $e(\ell_i)|\tilde B'|$, the remainder \cite[(1.16)]{Ba3} by $(e(\ell_i)^{1-\alpha})^{n+n'}$, the norm comparison between the norms \cite[(1.32)]{Ba3} at $B_0$ and at $\tilde B$ costing $(1+O(p_i))^{\#\mathrm{legs}}$, absorbed by the slack $\ell_i^{(4-d)/4-\beta}$ of the applications \cite[p.~421]{Ba3}); so a smaller cube only helps. (5) \emph{Gauge-out of $B_0$:} $B_0=\partial^\eta\lambda_0$ with $\lambda_0(x)=B_0\cdot(x-c)$ on $\square$; the Neumann covariant Laplacian on the bonds of $\square$ and the averaging $Q_i$ (contours inside blocks contained in $\square$) transform covariantly, $G_i(\square,B_0)=U(\lambda_0)G_i(\square,0)U(\lambda_0)^*$, and every class is invariant when the external scalar fields are transformed simultaneously \cite[(2.69)--(2.74)]{Ba2}; exact as long as $\square$ does not wrap around the torus (no condition on the cube). (6) \emph{$\Omega$-freeness after the split:} every old vertex lies in $\square_1\subset\square\subset\Omega$ and every new vertex in $\square$, so the characteristic function of $\Omega$ is $1$ wherever a vertex or a line end occurs; every object depending on $\Omega$ or on the cube ($\delta G_i(\Omega,\Omega_2,\tilde B)$, \cite[(1.16)]{Ba3}, the two $\delta G$'s, the gauge factors) has been moved into external fields, used on \cite[pp.~434--445]{Ba3} only through the norms \cite[(3.3)]{Ba3} at points of $\square_1$; the $\Omega_1$-dependent objects $g_i$, $1_{\Omega_1}$, $\dm(e',g_i,\lambda',\Omega_1,x)$ stay in the vertices, fixed before the step and entering only identities between a graph and its counterterm with the same restriction; on \cite[pp.~434--445]{Ba3} every propagator is $G_i(0)$, a slice, $G^\xi$, $C^\xi$ or $\Pi$, every sum runs over $\eta\Z^d$, $\xi\Z^d$, the cubes $\Delta(v)$, $\Omega_1$ or the torus identified with $\eta\Z^d$, Taylor's formula \cite[(3.10)]{Ba3} is applied to legs and to $g,g'$, and the $\lambda$-vertices are matched against counterterms with the same localisation \cite[(3.19), pp.~444--445]{Ba3}: none of this refers to the cube side or to $\dist(\cdot,\Omega^c)$. (7) \emph{Chains without old vertices} \cite[p.~433]{Ba3}: a class with no old vertex is a path $v_1$--$\cdots$--$v_m$ of new vertices of one product $(G_0V)^mG_0$, $G_0\defeq G_i(\square,B_0)$, with the cut lines at its ends as external scalar fields; the printed route through \cite[Prop.~2.1]{Ba3} is not used (its summation by parts \cite[(2.8)]{Ba3} is printed on the torus and on the Neumann box produces a face term where the $\tilde B'$-localisation does not vanish); instead, with $\norm\psi_E^2\defeq\ip\psi{G_0^{-1}\psi}$, $G_0^{-1}\ge D_{B_0}^*D_{B_0}$ and $G_0^{-1}\ge\gamma_\square$ (constant field, gauged out), each vertex of \cite[(1.8)--(1.11), (1.13)--(1.15)]{Ba3} is a bounded form in $\norm\cdot_E$: the form of \cite[(1.8)]{Ba3} is $\kappa\sum_{b\subset\square}\eta^d(D_{B_0}\phi)(b)\cdot M(b)\chi(b_-)$ with $|M(b)|\le C(\sup|\tilde B'|)^{n'}$, bounded through $\norm{D_{B_0}\phi}\le\norm\phi_E$ and $\norm\chi\le\gamma_\square^{-1/2}\norm\chi_E$; \cite[(1.9), (1.11)]{Ba3} likewise, with $|R_{\bar n+1}|\le1$; \cite[(1.10)]{Ba3} carries no derivative and uses $\gamma_\square^{-1/2}$ on both legs; a block pair \cite[(1.13)--(1.14)]{Ba3} uses $|\tilde B'(\Gamma)|\le2d\sup|\tilde B'|$ on its unit-block contours and $\sum_{x\in B(y)}\eta^d|\phi(x)|\le\norm\phi$; every form norm is $\le C(e(\ell_i)\sup|\tilde B'|)^{n'_v}$. Hence $E(H)=\ip{G_0^{1/2}K_1^*\Phi_1}{\hat K_2\cdots\hat K_{m-1}G_0^{1/2}K_m\Phi_m}$, $\hat K_v\defeq G_0^{1/2}K_vG_0^{1/2}$, is bounded by the product of the form norms times the end norms $\norm{h\Phi}_E$, where $h=1$ on the $2\eta$-neighbourhood of the end cube, $|\partial^\eta h|\le c_h$, and $\norm{h\Phi}_E^2=\norm{D_{B_0}(h\Phi)}^2+m^2\ell_i^2\norm{h\Phi}^2+a_i\norm{Q_i(B_0)h\Phi}^2\le C(1+c_h)^2\norm{h\Phi}_{1,0}^2$ by the Leibniz rule on a cube of side $\le4$; so $\norm{1_{\square^\eta(v)}\Phi}+\norm{1_{\square^\eta(v)}D_{B_0}\Phi}\le C\norm{h\Phi}_{1,0}$ is all that is used of an end field. The cut lines are bounded pointwise on $\square$, with their first differences in each slot, at pairs of boxes at distance $\ge1-2\eta\ge\frac35$, by \cref{lem:a01:RBK}: $|D^a_{B_0,x}D^b_{B_0,x'}G_0(x,x')|\le Ce^{-\delta|x-x'|}$ for $|x-x'|\ge\frac35$ and $a,b\in\{0,1\}$, at every pair of points of $\square$, its faces included, with no depth condition; and a cut line $K$ contracted between two functionals $A$, $B$ bounded by the sup-and-difference norms of the two end regions obeys $|A(x\mapsto B(K(x,\cdot)))|\le\norm A\norm B\max_{a,b\in\{0,1\}}\sup|D^a_xD^b_{x'}K(x,x')|$, because $x\mapsto B(K(x,\cdot))$ is linear in $K$ and its differences are $B(D_xK(x,\cdot))$ bond by bond, so no hypothesis $\dist\ge1$ on the two supports is needed and distance exactly $1$ is covered. A cut line between a chain end and an old-vertex class is an external scalar field of that class, read through its norms \cite[(3.3)]{Ba3} after localisation to the $\eta$-neighbourhood of the class's cube by a function $h'$ as in \cref{lem:a01:RBK}(c), the chain slot anywhere in $\square$: by the same linearity, with $B$ the class functional, bounded by Ba{\l}aban's norm \cite[(1.32)]{Ba3} of its external field, $|A(x\mapsto B(h'K(x,\cdot)))|\le\norm A\norm B\max_{a\in\{0,1\}}\sup_x\norm{h'D^a_xK(x,\cdot)}$ with the norm \cite[(1.32)]{Ba3} in the last factor, and \cref{lem:a01:RBK}(c) bounds that factor by $C(1+c_h)^2$ at every $x$ of the chain-end region, the faces of $\square$ included. The vertex forms are sums over the bonds of $\square$, so vertices at the faces need no separate treatment. (8) \emph{The small factors of \cite{Ba2} survive:} the factors $O(e^{-\delta_1r_i})$ that \cite[p.~422]{Ba3} produces at domain changes come from $\exp[-\delta_0d(\{\square(v)\})]$ in \cite[(1.33)]{Ba3}, i.e.\ from the lines removed before the cube step, whose bound depends on $\Omega$ only; the two-domain inputs enter as external fields, linearly. So the only place where $r_\square=r_i$ was used is (1), where $r^\flat$ serves by hypothesis; the constant acquires the factors of (3) and (4), both of the printed kind. For the last claim: a vertex anchored at $\Omega_1$ has depth $\ge5r(\ell_{i-1})/5-3\ge r_i-3\ge R-3\ge D^\flat$ (\cref{lem:a01:X0}(a), $r(\ell_{i-1})\ge r_i\ge R$); a vertex anchored at $\supp\tilde A$ has depth $>2r_i-2-M_\theta-2\eta\ge2R-2.4-M_\theta>D^\flat$ (\cref{lem:a01:X0}(a)); the other steps of the proof need depth at most $R_{\mathrm{Ba}}+13<D^\flat$ at old vertices or within $2\sqrt3+1$ of $\supp\tilde A$.
\end{proof}
% a01_H0rel part 6x1, continued: two-slot bounds of the box propagator at a constant field (input from a01_p6x1.tex after prop:a01:P12)
\begin{lemma}[two-slot bounds of the box propagator at a constant field]\label{lem:a01:RBK}
Let $i\ge1$, $\eta\defeq L^{-i}$ the fine spacing in scale-$i$ units, $\square$ a box of large blocks (a product of intervals, each a union of large blocks) that does not wrap around the torus, with Neumann conditions on its faces, $B_0$ a constant field, and $G_0\defeq G_i(\square,B_0)$ the propagator of step (7) of the proof of \cref{prop:a01:P12}, $G_0^{-1}=D_{B_0}^*D_{B_0}+m^2\ell_i^2+a_iQ_i(B_0)^*Q_i(B_0)$ on $\ell^2(\square\cap\eta\Z^d)$ with the uniform block average $Q_i$. Kernels are taken for the measure $\eta^d$, and $D^a_{B_0,x}$, $a\in\{0,1\}$, denotes the identity or a covariant forward difference on a bond of $\square$ at $x$. \textup{(a)} There are $C,\delta>0$ depending on $(d,L,a)$ and the mass range only, free of $i$, $\eta$, $\square$, $B_0$ and the positions, such that for all $x,x'\in\square$ and $a,b\in\{0,1\}$ ($d\ge3$; here $d=3$)
\[
|D^a_{B_0,x}D^b_{B_0,x'}G_0(x,x')|\ \le\ C\,(\eta+|x-x'|)^{-(d-2+a+b)}\,e^{-\delta|x-x'|},
\qquad\text{so}\qquad |D^a_{B_0,x}D^b_{B_0,x'}G_0(x,x')|\le5C\,e^{-\delta|x-x'|}\ \text{ for }|x-x'|\ge\tfrac35\ (d=3).
\]
\textup{(b)} No depth condition enters: the bound holds at every pair of points of $\square$, the bonds of its faces included.

\textup{(c)} \emph{One slot in Ba{\l}aban's norm.} Let $\Delta,\Delta'$ be two cubes of $\square$ at distance $\ge1$, let $h'$ be a localisation function on $\square$, $0\le h'\le1$, equal to $1$ on the $\eta$-neighbourhood of $\Delta'$ and vanishing at distance $\ge\frac12$ from $\Delta'$, with $|\partial^\eta h'|\le c_h$ and $|\partial^\eta\partial^\eta h'|\le c_h$, and let $x$ lie within $\eta$ of $\Delta$. Then for $a\in\{0,1\}$ the field $\Psi\defeq h'(\cdot)\,D^a_{B_0,x}G_0(x,\cdot)$ satisfies the three clauses of Ba{\l}aban's external-field norm \cite[(1.32)]{Ba3} (sup, first differences, $\alpha_0$-H\"older quotients of the first differences on unit balls) on all of $\square$, its faces included, with the bound $C(1+c_h)^2$, $C$ of the same kind as above.
\end{lemma}
\begin{proof}
(1) \emph{Gauge-out.} With $\lambda_0(x)\defeq B_0\cdot(x-c)$ on the non-wrapping $\square$, $G_i(\square,B_0)=U(\lambda_0)G_i(\square,0)U(\lambda_0)^*$ and $D_{B_0}(U(\lambda_0)\phi)=U(\lambda_0)D\phi$ bond by bond (step (5) of the proof of \cref{prop:a01:P12}; $|U|=1$), so every modulus in the claim equals the one at $B_0=0$. Let $B_0=0$.

(2) \emph{Even images.} Let $\square=\prod_\mu[\alpha_\mu,\beta_\mu)$ and let $T_\square$ be the torus of periods $2(\beta_\mu-\alpha_\mu)$, each a multiple of $2s_b$, so that the large blocks of $\square$ tile $T_\square$ periodically; let $\mathcal W_\square\cong(\Z_2)^d$ be the group generated by the reflections in the face planes $x_\mu=\alpha_\mu-\eta/2$ of $\square$ (planes between sites on block boundaries; on $T_\square$ the reflection in $x_\mu=\beta_\mu-\eta/2$ is the same element). They map sites to sites, bonds to bonds and blocks to blocks, the uniform average commutes with them, and the mass term is constant; every bond of $T_\square$ that crosses a face plane joins a site to its mirror image, where a $\mathcal W_\square$-even field has zero difference, and the remaining bonds of $T_\square$ inside $\square$ are the bonds of $\square$. Hence, exactly as in \cref{lem:a01:PE}(a) and \cref{lem:a01:NL}, the Neumann operator on $\square$ is the even restriction of the periodic one, $G_i(T_\square,0)\Ebar=\Ebar\,G_i(\square,0)$ with $\Ebar$ the even extension, and
\[
G_i(\square,0)(x,x')=\sum_{g\in\mathcal W_\square}G_i(T_\square,0)(x,gx')\qquad(x,x'\in\square),
\]
a difference in $x'$ along a bond of $\square$ acting on the image terms along the reflected bonds. Every image satisfies $\dist_{T_\square}(x,gx')\ge|x-x'|$: coordinatewise, for $x_\mu,x'_\mu\in[\alpha_\mu,\beta_\mu-\eta]$ and $r$ the reflection in $\beta_\mu-\eta/2$, $rx'_\mu-x_\mu=|x_\mu-x'_\mu|+2(\beta_\mu-\tfrac\eta2-\max(x_\mu,x'_\mu))\ge|x_\mu-x'_\mu|$, and the complementary distance around the period is $x_\mu+x'_\mu-2\alpha_\mu+\eta\ge|x_\mu-x'_\mu|$.

(3) \emph{Periodisation.} At zero field $G_i(T_\square,0)(x,y)=\sum_{n\in\Z^d}G_i(\eta\Z^d,0)(x,y+2\ell n)$, $\ell\defeq(\beta_\mu-\alpha_\mu)_\mu$ (the periodic operator is the periodisation of the free-lattice one, as in step (3) of the proof of \cref{prop:a01:P12}), every term having $|x-y-2\ell n|\ge\dist_{T_\square}(x,y)$, and at most $(2+k)^d$ values of $n$ have $|x-y-2\ell n|<k+1$, every period being $\ge1$. So, with the free-lattice bound (4) below at the rate $\delta'$ and $\rho\defeq|x-x'|$,
\[
|D^a_xD^b_{x'}G_i(\square,0)(x,x')|\ \le\ 2^dC'(\eta+\rho)^{-(d-2+a+b)}\sum_{k\ge\lfloor\rho\rfloor}(2+k)^de^{-\delta'k}\ \le\ C(\eta+\rho)^{-(d-2+a+b)}e^{-\delta'\rho/2}.
\]

(4) \emph{The free lattice.} On $\eta\Z^d$ at zero field, $G_i(\eta\Z^d,0)=G_1+\sum_{j=1}^{i-1}G_{(j)}$ with the slices $G_{(j)}\defeq G_{j+1}-G_j=a_jG_jQ^*C^{(j)}QG_j$ (the Schur identity behind \cref{lem:a01:legs}(b) and \cite[(2.17)]{K1}, with $C^{(j)}$ as in \cref{lem:a01:A23}(c)), so that the sum telescopes to $G_i$. In unit-$j$ units ($j$-blocks of side one, spacing $L^{-j}$) let $\ell_j(\cdot,u)\defeq a_jG_jQ^*e_u$ be the leg of the unit datum on the $j$-block $u$; since $G_j$ is real symmetric at zero field, $G_{(j)}(y,y')=a_j^{-1}\sum_{u,v}\ell_j(y,u)C^{(j)}(u,v)\ell_j(y',v)$, each variable sitting in its own leg, so
\[
D^a_yD^b_{y'}G_{(j)}(y,y')=a_j^{-1}\sum_{u,v}(D^a\ell_j)(y,u)\,C^{(j)}(u,v)\,(D^b\ell_j)(y',v).
\]
The legs obey $|D^a\ell_j(y,u)|\le Ce^{-\delta_2\dist(y,B(u))}$ at every site and every bond of the free lattice, $a\in\{0,1\}$ (the leg bounds \cite[(3.64)]{K1}, \cite[(2.12)]{Ba3}, at the uniform average $f=1$ in \cref{prop:a02:FL}; on the free lattice every bond is interior, so they are also \cref{lem:a01:A4}(b) and the interior case (4) of the proof of \cref{lem:a01:A5}), and $|C^{(j)}(u,v)|\le Ce^{-\delta_1|u-v|}$ (\cite[Thm.~3.3, (3.6)]{K1}; \cref{lem:a01:A23}(c)). With $|y-y'|\le\dist(y,B(u))+|u-v|+\dist(y',B(v))+\sqrt d$ and two geometric sums over $u,v\in\Z^d$, $|D^a_yD^b_{y'}G_{(j)}(y,y')|\le C_Se^{-\delta_S|y-y'|}$ at every pair of points, $\delta_S\defeq\frac12\min\{\delta_1,\delta_2\}$. In scale-$i$ units the slice $j$ has the kernel $s_j^{2-d}G_{(j)}(x/s_j,x'/s_j)$ with $s_j\defeq L^{j-i}$, and a scale-$i$ difference is $s_j^{-1}$ times a unit-$j$ one, so slice $j$ contributes at most $C_Ss_j^{-(d-2+a+b)}e^{-\delta_S|x-x'|/s_j}$; $G_1$, read in unit-$1$ units at the fixed spacing $L^{-1}$, obeys the same bound with $s\defeq L\eta$ and a constant depending on $L$ (its kernel is $L^d\ip{e_y}{G_1e_{y'}}$, bounded with the decay of \cref{lem:a01:A23}(a), and a difference at spacing $L^{-1}$ costs a factor $2L$). Finally, for $p\defeq d-2+a+b>0$ and $s$ ranging over the powers of $L$ in $[L\eta,L^{-1}]$: $e^{-\delta_S\rho/s}\le e^{-\delta_S\rho/2}e^{-\delta_S\rho/(2s)}$ as $s\le1$; the terms with $s\ge\max(\rho,\eta)$ sum to at most $(1-L^{-p})^{-1}\max(\rho,\eta)^{-p}$, and the terms with $s<\rho$ are $\rho^{-p}(\rho/s)^pe^{-\delta_S\rho/(2s)}$ with $\rho/s$ running over powers of $L$, summable to at most $C_p\rho^{-p}$. Hence $|D^a_xD^b_{x'}G_i(\eta\Z^d,0)(x,x')|\le C'(\eta+\rho)^{-p}e^{-\delta'\rho}$ with $\delta'\defeq\delta_S/2$, which is (4), and (3) gives the claim with $\delta\defeq\delta'/2$. The last clause of the statement is $(\eta+\rho)^{-p}\le(5/3)^3<5$ for $\rho\ge\frac35$.

(5) \emph{The clause (c).} By (1) every modulus is the one at $B_0=0$. The reflections of (2) and the translations of (3) are isometries of the lattice mapping bonds to bonds, so each clause of $x'\mapsto D^a_xG_i(\square,0)(x,x')$ at a pair $x',y'\in\square$ is bounded by the sum over the images of the same clause of $x'\mapsto D^a_xG_i(\eta\Z^d,0)(x,x')$ at the pair $(gx'+2\ell n,gy'+2\ell n)$, whose points are at least as far from $x$ as $x'$, $y'$ are, with the count of (3). Write $\Psi_0(x')\defeq D^a_xG_i(\eta\Z^d,0)(x,x')$ and $\rho'\defeq\min\{|x'-x|,|y'-x|\}$. \emph{Claim:} on the free lattice at zero field the sup and first-difference clauses of $\Psi_0$ at $x'$, and the H\"older quotient at a pair $x',y'$ with $|x'-y'|\le1$, are $\le Ce^{-\delta\rho'}$ whenever $\rho'\ge\frac1{20}$. Indeed, by the factorisation of (4), in unit-$j$ units $D^a_xG_{(j)}(x,\cdot)=\sum_v\ell_j(\cdot,v)\,a_j^{-1}\beta_j(v)$ is the leg of the block data $a_j^{-1}\beta_j$, $\beta_j(v)\defeq\sum_u(D^a\ell_j)(x,u)C^{(j)}(u,v)$, with $|\beta_j(v)|\le Ce^{-\delta_S\dist(x,B(v))}$ by the leg and $C^{(j)}$ bounds of (4) and one geometric sum. Split the data into the blocks with $\dist(x,B(v))\le\rho'_j/2$, $\rho'_j\defeq\rho'/s_j$ the distance in unit-$j$ units, and the rest. The leg of bounded block data obeys the sup bound of \cref{lem:a01:A4}(b), the first-difference bound of \cref{lem:a01:A5} and its H\"older bound on balls of radius $\frac1{16}$ with decay $e^{-\delta_3\dist(\cdot,S)}$ from the support $S$ of the data (the interior case, every bond of the free lattice being interior; printed: \cite[(3.63)--(3.65)]{K1}, \cite[(2.12)]{Ba3}); the far part has data of size $\le Ce^{-\delta_S\rho'_j/2}$, and the near part has its support within $\rho'_j/2$ of $x$, so at distance $\ge\rho'_j/2$ from $x'$ and $y'$; a pair at unit-$j$ distance $>\frac1{16}$ is bounded by $2\cdot16^{\alpha_0}$ times the first differences. So in unit-$j$ units every clause of slice $j$ is $\le Ce^{-\delta_6\rho'_j}$, $\delta_6\defeq\frac14\min\{\delta_S,\delta_3\}$. In scale-$i$ units the sup clause of slice $j$ scales as $s_j^{-(d-2+a)}$, the first differences as $s_j^{-(d-1+a)}$ and the H\"older quotients as $s_j^{-(d-1+a+\alpha_0)}$ (a scale-$i$ pair at distance $t\le1$ is a unit-$j$ pair at distance $t/s_j$; for $t>s_j/16$ use two first differences, $2\cdot16^{\alpha_0}s_j^{-(d-1+a)}t^{-\alpha_0}\le2\cdot16^{\alpha_0}s_j^{-(d-1+a+\alpha_0)}$); the slice $G_1$ at the fixed spacing contributes a bounded multiple of its first differences to the H\"older clause, every pair being at distance $\ge\eta$. The sum over $s$ is the geometric sum of (4) with $d-1+a+\alpha_0$ in place of $p$, giving $\le C(\eta+\rho')^{-(d-1+a+\alpha_0)}e^{-\delta'\rho'}\le Ce^{-\delta'\rho'}$ for $\rho'\ge\frac1{20}$; this is the claim, and with the image count it holds for $x'\mapsto D^a_xG_i(\square,0)(x,x')$ at pairs of $\square$ with $\rho'\ge\frac1{20}$. \emph{Localisation.} $\Psi=h'\Psi_\square$, $\Psi_\square\defeq D^a_xG_i(\square,0)(x,\cdot)$, and $D(h'\Psi_\square)=h'(\cdot+\eta e)D\Psi_\square+(\partial^\eta h')\Psi_\square$ bond by bond. Both $\Psi$ and $D\Psi$ vanish outside $S'\defeq\{\dist(\cdot,\Delta')\le\frac12+\eta\}$, whose points are at distance $\ge\frac12-2\eta\ge\frac1{10}$ from $x$ ($\eta\le\frac15$, $x$ within $\eta$ of $\Delta$, $\dist(\Delta,\Delta')\ge1$). Hence the sup and first-difference clauses of $\Psi$ are those of $\Psi_\square$ on $S'$ times $1+c_h$. For a H\"older pair $x',y'$ with both points at distance $\ge\frac1{20}$ from $x$, the Leibniz rule and the bounds $|h'(x')-h'(y')|\le dc_h|x'-y'|$, $|\partial^\eta h'(x')-\partial^\eta h'(y')|\le dc_h|x'-y'|$ give the quotient $\le C(1+c_h)^2$ from the three clauses of $\Psi_\square$; a pair with one point $y'$ within $\frac1{20}$ of $x$ has $D\Psi(y')=0$ and $|x'-y'|\ge\frac1{10}-\frac1{20}=\frac1{20}$ whenever $D\Psi(x')\ne0$, so its quotient is $\le20^{\alpha_0}$ times a first difference. This is (c).
\end{proof}

On the cube the same device, a constant box $Q_x$ of large blocks around each vertex cube at depth $\ge R_*$, is step (1) of the proof of \cref{thm:a01:B2}; the depth supply is \cref{lem:a01:X0}(a) there too.

% a01_H0rel part 6x2: the replaced steps (b)-(d) at block factor 5 (input from a01_H0rel_p6.tex after a01_p6x1)
\subparagraph{(b) The cut-off $\zeta^{(i)}$.} In \cite[(2.44)]{Ba2} $A^{(i)}=a_i\ell_i^{-2}\zeta^{(i)}G_iQ_i^*A$ with $\supp\zeta^{(i)}(\cdot,y)\subset\{|x-y|<r_i-2M\}$, $\zeta^{(i)}(x,y)=1$ for $|x-y|\le\frac12r_i$. The printed proof of \cite[(2.60)]{Ba2}, $\partial^\eta A^{(i)}=O(p_i)$ on $\Omega$, bounds the first two terms of \cite[(2.61)]{Ba2} by the restrictions \cite[(2.55)]{Ba2} at every coarse point $y'$ with $\zeta^{(i)}(x,y')\ne0$; \cite[(2.55)]{Ba2} is available on $\Lambda_{-1}^{(i-1)\prime}$, which contains $\Omega$ at depth $\ge3r(\ell_{i-1})/L-2$ (\cref{lem:a01:X0}(a), three layers). The containment $3r(\ell_{i-1})/L-2\ge r_i-2M$ has $\alpha=1$, $\beta=3$: it holds at every scale iff $L\le3$. The gap is real: by \cite[(2.9)]{Ba2} the restrictions survive the partial resummation only on $\Lambda'_{-1}$, $\Lambda^*_{-1}$, $\Lambda_{-1}$, so at a large block which stayed large at every finer scale $A(y')$ is unrestricted and the contribution $a_i\zeta^{(i)}(x,y')(G_iQ_i^*)(x,y')A(y')$ has no bound; hence \cite[(2.60)]{Ba2} as printed is not established in a collar of width $\approx r_i-3r(\ell_{i-1})/L$ along $\partial\Omega$, while \cite[p.~571]{Ba2} uses it on all of $\Omega$ and \cite[Prop.~1]{Ba3} assumes it on \cite[p.~412]{Ba3}. The same cut-off enters the split \cite[(2.89)]{Ba2} at the next scale.

\begin{proposition}[the shrunk cut-off]\label{prop:a01:Zth}
Replace $\zeta^{(i)}$ in \cite[(2.44)]{Ba2}, at every scale, by $\zeta^{(i)}_\theta$ with $\supp\zeta^{(i)}_\theta(\cdot,y)\subset\{|x-y|<\theta r_i-2M\}$, $\zeta^{(i)}_\theta(x,y)=1$ for $|x-y|\le\frac12\theta r_i$, $|\partial^\eta\zeta_\theta|\le1$, $|\partial^\eta\partial^\eta\zeta_\theta|\le2$ (a $C^{1,1}$ radial profile), with $\theta=\theta_5\defeq\frac25$; the first-step cut-off of \cite[p.~561]{Ba2} already has range $r/(2L)$ in the units of the next scale, so this returns to the first-step convention. Let $R\ge\max\{10(M+1),\,M_\theta+2,\,(M_\theta+7)/11.6\}$, where $\theta_{i+1}$ equals $1$ on $B^i(\Lambda_2^{(i)})$ and $0$ at distance $\ge M_\theta$ from it ($M_\theta<M$ by \cite[p.~567]{Ba2}, where $\theta_{i+1}$ ``varies smoothly'' on a slice of thickness $<M$; we take $M_\theta=2M/3$ and a $C^{1,1}$ profile). Then at $L=5$:
\begin{enumerate}[label=\textup{(Z\arabic*)},nosep,leftmargin=*]
\item for $y\in\Lambda_2^{(i-1)\prime}$, $x\in B^i(y)$, the $\zeta_\theta$-ball $\{|x-y'|<\theta r_i-2M\}$ lies in $\Lambda_{-1}^{(i-1)\prime}$, where \cite[(2.55)]{Ba2} holds;
\item \cite[(2.59)--(2.60), (2.64)]{Ba2} hold with $O(p_i)$ replaced by $O_L(p_i)$, the constant containing $p_{i-1}/p_i\le(1+\ln L)^p$;
\item the terms carrying $1-\zeta^{(i)}_\theta$ in \cite[(2.61), (2.64)]{Ba2}, the terms carrying $1-\zeta^{(i)}_\theta$, $1-\zeta^{(i+1)}_\theta$ in \cite[(2.88)--(2.90)]{Ba2}, and the second term of \cite[(2.89)]{Ba2} (the one with $\delta C^{(i)}_{\Lambda_0}$) are $O(\ell_i^\kappa)$ for every $\kappa$, with their derivatives and the $\alpha_0$-H\"older quotients of their derivatives;
\item the transition width $\frac12\theta r_i-2M\ge2$ required by the $C^{1,1}$ radial profile with $|\partial^\eta\zeta_\theta|\le1$ and $|\partial^\eta\partial^\eta\zeta_\theta|\le2$ holds;
\item the split \cite[(2.89)]{Ba2} holds \emph{as printed}, with the full $B$ in its fifth term, and the $\zeta^{(l)}_\theta$-balls around $\supp\theta_l$ lie in $\Lambda_{-1}^{(l-1)\prime}$ at every scale $l$.
\end{enumerate}
No $\zeta$ obeying the printed \cite[(2.44)]{Ba2} makes \cite[(2.89)]{Ba2} exact at $L=5$ for small $\ell_i$ (the minimal printed support $\frac12r$ gives $L\theta=\frac52>2$), so the shrunk cut-off is necessary at $L=5$ if \cite[(2.89)]{Ba2} is kept in Ba{\l}aban's form.
\end{proposition}
\begin{proof}
(Z1): $\theta r_i\le3r(\ell_{i-1})/L-2+2M$ is \cref{lem:a01:X0}(b) with $\alpha=\frac25$, $\beta=3$, $\alpha L=2\le3$, margin $\ge r(\ell_{i-1})/5-2\ge0$ for $R\ge10$. (Z2): in \cite[(2.61)]{Ba2} with $\zeta_\theta$ the first term sums over $y'$ in the $\zeta_\theta$-ball, where by (Z1) $|A(y')-A(y)|\le c_1p_{i-1}L|y'-y|$ (a path of $\le L|y'-y|$ bonds with $|\partial A|\le c_1p_{i-1}$) and \cite[Prop.~2.2]{Ba2} gives the kernel bound $c_{\mathrm{Ba}}e^{-\delta_0|x-y'|}$, so the term is $O(p_{i-1})=O_L(p_i)$; the third term is the exact \cite[(2.62)--(2.63)]{Ba2} ($G_iQ_i^*1$ is constant); the second is a tail of (Z3); \cite[(2.64)]{Ba2} is the same for the derivative, with $\partial^\eta(G_iQ_i^*1)=0$. (Z3): the $(1-\zeta_\theta)$ term of \cite[(2.61)]{Ba2} multiplies the single constant $A(y)$, $|A(y)|\le c_1p_{i-1}/(\lmu\ell_{i-1})$, with a kernel supported at $|x-y'|\ge\frac12\theta r_i$ and bounded by $c_{\mathrm{Ba}}e^{-\delta_0|x-y'|}$, so it is $\le Cr_i^de^{-\delta_0\theta r_i/2}p_{i-1}/(\lmu\ell_{i-1})\le\ell_i^\kappa$ for small $\ell_i$ (\cref{lem:a01:X0}(c)); in \cite[(2.89)]{Ba2} the third and fourth terms carry $1-\zeta^{(i)}_\theta$, $1-\zeta^{(i+1)}_\theta$ acting on $\Lambda_0^{(i)\prime}B$ with $|B|\le c_1p/(\lmu\ell)$ there, and the same bound gives them, their derivatives and H\"older quotients (kernel by kernel: \cite[(2.34), (2.36)]{Ba1} for $C^{(i)}$, $\delta C_{\Lambda_0}$, the composition identity of (Z5) for $G^\eta_{i+1}Q_{i+1}^*$, the $C^{1,1}$ profile at scales $i$, $i+1$); the second term $T_2(x)=a_iaL^{-2}\theta_{i+1}(x)\sum_y\zeta^{(i)}(x,y)(G_iQ_i^*)(x,y)(\delta C^{(i)}_{\Lambda_0}Q^*\Lambda_0^{(i)\prime}B)(y)$ has, for $x\in\supp\theta_{i+1}$ and $\zeta(x,y)\ne0$, $\dist(y,\Lambda_0^{(i)c})>r_i-M_\theta-2$ (two layers of \cite[(2.8)]{Ba2} minus the range of $\zeta$), and $\delta C_{\Lambda_0}(y,y')$ decays exponentially in $|y-y'|$ and in the distances of $y$, $y'$ to $\Lambda_0^{(i)c}$ \cite[(2.35)]{Ba1}, so $|T_2|+|\partial^\eta T_2|\le Cr_i^de^{-\delta'(r_i-M_\theta-2)}p/(\lmu\ell)\le C_\kappa\ell_i^\kappa$. (Z4): $\frac12\theta r_i-2M\ge2$ iff $r_i\ge(4+4M)/\theta=10(M+1)$, and $r_i\ge R$. (Z5): on the unit lattice of scale $i$, with $Q$ the one-step average, $P_i=Q_i^*Q_i$, $P_{i+1}=Q_i^*Q^*QQ_i$, $G_i=(H+a_iP_i)^{-1}$, $G^\eta_{i+1}=(H+a_{i+1}L^{-2}P_{i+1})^{-1}$, $\Delta^{(i)}=a_i-a_i^2Q_iG_iQ_i^*$, $C^{(i)}=(aL^{-2}Q^*Q+\Delta^{(i)})^{-1}$ and $a_{i+1}=aa_i/(a_i+aL^{-2})$ \cite[(2.13), (2.21), (2.31)]{Ba1}, one has the composition identity $a_iaL^{-2}G_iQ_i^*C^{(i)}Q^*=a_{i+1}L^{-2}G^\eta_{i+1}Q^*_{i+1}$: apply $H+a_{i+1}L^{-2}P_{i+1}$ to the left side, use $HG_i=1-a_iP_iG_i$, $1-a_iQ_iG_iQ_i^*=\Delta^{(i)}/a_i$, $\Delta^{(i)}C^{(i)}=1-aL^{-2}Q^*QC^{(i)}$ and $Q^*QQ^*=Q^*$; the coefficient of $Q^*QC^{(i)}Q^*$ vanishes iff $a_{i+1}(a_i+aL^{-2})=aa_i$, which is \cite[(2.13)]{Ba1}, and the remaining coefficient is $a_{i+1}L^{-2}$ by the same relation (this is \cite[(2.41)]{Ba1} in adjoint, rescaled form). By \cite[(2.80)]{Ba2} and the Dirichlet condition \cite[(2.32)]{Ba1}, $C^{(i)}_{\Lambda_0}Q^*B=C^{(i)}_{\Lambda_0}Q^*\Lambda_0^{(i)\prime}B$; writing $C_{\Lambda_0}=C^{(i)}+\delta C_{\Lambda_0}$ and $\zeta=1-(1-\zeta)$ and using the identity, $\theta_{i+1}A^{(i)}$ equals the five terms of \cite[(2.89)]{Ba2} with $\Lambda_0^{(i)\prime}B$ in the fifth term, identically in $A'$, $B$; the printed form differs by the defect $a_{i+1}L^{-2}\theta_{i+1}\zeta^{(i+1)}_\theta G^\eta_{i+1}Q^*_{i+1}\Lambda_0^{(i)\prime c}B$, which vanishes when $\zeta^{(i+1)}_\theta(x,y)=0$ for all $x$ with $\theta_{i+1}(x)\ne0$ and $y\notin\Lambda_0^{(i)\prime}$, i.e.\ when $L(\theta r_{i+1}-2M)\le2r_i-M_\theta-2-L$ (such $x$ lie within $M_\theta+1$ of $\Lambda_2^{(i)}$, such $y$ within $L$ of $\Lambda_0^{(i)c}$, and $\dist(\Lambda_2^{(i)},\Lambda_0^{(i)c})>2r_i-1$). With $L\theta_5=2$ this reads $2(r_i-r_{i+1})\ge M_\theta+7-10M$, and $r_i-r_{i+1}=R\ln5\,(2u_{i+1}+\ln5)\ge5.81R$. The balls around $\supp\theta_l$ reach at most $2r_l-10M+M_\theta+1$ beyond $\Lambda_2^{(l-1)}$ (scale-$(l-1)$ units), while $\Lambda_{-1}^{(l-1)}$ contains the $(3r(\ell_{l-1})-2)$-neighbourhood of $\Lambda_2^{(l-1)}$ (\cref{lem:a01:X0}(a)): $2r_l+M_\theta+8\le3r(\ell_{l-1})+10M$ holds since $r_l\le r(\ell_{l-1})$ and $r(\ell_{l-1})\ge R\ge M_\theta+2$. For the printed $\theta=1$ the defect condition has $L\theta=L$, $\beta=2$, so it holds at every scale iff $L\le2$; on the reading that absorbs the defect into $\tilde B^{(i)}$, the $\zeta^{(i+1)}$-balls must stay inside $\Lambda_{-1}^{(i)\prime}$, $\beta=3$, hence $L\le3$.
\end{proof}
The other uses of \cite[Lemma~2.3]{Ba2} are covered by (Z2): \cite[(2.68)]{Ba2} on the box of (c) below; \cite[(2.95)--(2.98)]{Ba2} at each scale $l$ together with (Z5); the regularity of $\tilde B$ needed by \cite[Prop.~1]{Ba3} at the background $\theta_iA^{(i)}$ on $\Omega$, where $\theta_i=1$; and the first-step use \cite[(2.25)]{Ba2}, whose cut-off has $\theta=1/(2L)<3/L$. On the consumer side, the all-small-field term of \cite[(3.4)]{K1} is affected only through the cut background, and the difference between the actions at $A^{(j)}$ and at $A_\theta$ is $O(\eta^{s'})|X^{(j)}|$ on $\supp\chi_j$, as used in the proof of \cref{prop:a01:upper}.

\subparagraph{(c) The boxes of \cite[Lemma~2.4]{Ba2}.} The printed proof \cite[p.~572]{Ba2} takes $\square_1,\square_2$ as the sums of large blocks of $\Lambda_7^{(i-1)\prime}$ within $2r_i$, $4r_i$ of $y$ and $\square=B^i(\square_2)$. For $y\in\Lambda_7$ at distance $O(1)$ from $\Lambda_7^c$ the step \cite[(2.67)]{Ba2} is not justified: the dropped source $\Lambda_6\phi-\square_1\phi$ contains $\phi$ on $\Lambda_6\setminus\Lambda_7$ at distance $O(1)$ from $x$, where $a_iG_iQ_i^*$ is $O(1)$, so this error is $O(\Upsilon_i)$, $\Upsilon_i\defeq c_1(\lambda\ell_{i-1})^{-1/4}p_{i-1}$, not $O(\ell_i^\kappa)$; the domain change $\Omega\to\square$ happens at distance $O(1)$ from $x$, so \cite[(2.26)]{Ba1} gives no smallness; and $x$ may lie within $R_{\mathrm{Ba}}$ of $\partial\square$, where \cite[(2.24)--(2.25)]{Ba1} need a parallelepiped with a constant field near its boundary, while $\Lambda_7\cap(\text{cube})$ has re-entrant edges. An untruncated box of radius $4r_i$ inside $\Omega$ would need $4r_i\le5\rho_i$, which fails for small $\ell_i$ at every $L\ge2$ ($\alpha=4$, $\beta=5$). The only consumer is the quartic convergence factor \cite[(2.78)--(2.79)]{Ba2}, a sum over $y\in\Lambda_7^{(i-1)\prime}\cap\Lambda_7^{(i)c}$, which contains every point of $\Lambda_7$ within $2r_i$ of $\Lambda_7^c$: the edge points are exactly the ones used.

\begin{lemma}[{\cite[Lemma~2.4]{Ba2}} with a complete proof, every $L$]\label{lem:a01:lem24}
Assume \cite[(2.55)]{Ba2}, and: $5(r(\eps_0^{\mathrm{Ba}})-s_b)/L\ge R_{\mathrm{Ba}}+c_d+2s_b$, $4r(\eps_0^{\mathrm{Ba}})\ge R_{\mathrm{Ba}}+c_d$, $R\ge10s_b+10$, with $c_d\defeq\sqrt d\,s_b$; $\eps_0^{\mathrm{Ba}}$ small enough that $r(\ell_{i-1})/r(\ell_i)\le2L$ for all $i$ (at $L=5$, $\varrho=2$ this is $(1+\ln5)^2\le10$, true for all $\eps_0^{\mathrm{Ba}}\le1$); the mass condition $\Upsilon_im^2\ell_i^2(G_i(\square,0)Q_i^*1)(x)=O(p_i)$ (Ba{\l}aban's own use of \cite[(2.75)]{Ba2}); and $Ca_i\Upsilon_i(1+r_i)^{d-1}\exp(-\delta_0[(r(\ell_{i-1})-s_b)/L-c_d])\le\ell_i^\kappa$, which holds for all small $\eps_0^{\mathrm{Ba}}$ at every $R>0$ because $\varrho>1$. Then \cite[(2.65)--(2.66)]{Ba2} hold for every $y\in\Lambda_7^{(i-1)\prime}$, $x\in B^i(y)$, $b\subset B^i(\Lambda_7^{(i-1)\prime})$, with constants independent of $\eps$, $i$, the lattice and the position of $y$ relative to $\partial\Lambda_7$.
\end{lemma}
\begin{proof}
Put $s_i\defeq\min\{4r_i,5\rho_i-2s_b\}$, let $\square_2$ be the cube of large blocks of sup-radius $s_i$ around the large block of $y$, $\square\defeq B^i(\square_2)$, and $\square_1\phi\defeq1_{Q(y,s_i/2)}\Lambda_6\phi$ with $Q(y,s)$ the sup-ball. Since $\dist(\Lambda_7,\Lambda_2^c)\ge5\rho_i$ (\cref{lem:a01:X0}(a)), $\square\subset\Omega$; $\square$ is a rectangular parallelepiped and a union of large blocks, and $s_i-c_d\ge R_{\mathrm{Ba}}$ at every $i$ by the hypotheses. (1) \emph{Source cut:} $\Lambda_6\phi-\square_1\phi$ is supported at distance $\ge s_i/2-c_d$ from $x$, and $\dist(x,\Omega^c)\ge5\rho_i-c_d\ge R_{\mathrm{Ba}}$, so \cite[(2.25)]{Ba1} on $\Omega$ with the sup bound $\Upsilon_i$ on $\phi$ gives $|a_i(G_i(\Omega,A^{(i)})Q_i^*(\Lambda_6\phi-\square_1\phi))(x)|\le Ca_ie^{-\delta_0(s_i/2-c_d)}\Upsilon_i$. (2) \emph{Domain change:} $G_i(\square,A)=G_i(\Omega,A)+\delta G_i(\square,\Omega,A)$, and by \cite[(2.26)]{Ba1} the second term applied to $Q_i^*\square_1\phi$ at $x$ carries $e^{-\delta_0(3s_i/2-2c_d)}$; the same for $D^\eta$ by \cite[(2.58)]{Ba2}. Hence $\phi^{(i)}(x)=a_i(G_i(\square,A^{(i)})Q_i^*(A^{(i)})\square_1\phi)(x)+O(a_i\Upsilon_ie^{-\delta_0(s_i/2-c_d)})$, and likewise for the covariant derivative: this is \cite[(2.67)]{Ba2} with the boxes chosen above. (3) \emph{Constant background:} $A_0\defeq A^{(i)}(y)$; joining points of the convex box by straight lattice paths and using \cite[(2.60)]{Ba2} (\cref{prop:a01:Zth}(Z2)) on $B^i(\Lambda_2^{(i-1)\prime})\supset\square$, $|A^{(i)}(z)-A_0|\le O(p)\sqrt d\,s_i\le O(pr_i)$ on $\square$; the expansion \cite[(3.44)]{Ba1} and \cite[Prop.~2.2]{Ba1} on the parallelepiped $\square$ replace $A^{(i)}$ by $A_0$ with error $O(\ell_i^{\kappa_0})$, as in \cite[(2.68)]{Ba2}. (4) \emph{Gauge-out:} $A_0=\partial^\eta\lambda_0$ with $\lambda_0(z)=A_0\cdot(z-y)$ on $\square$, so $U(A_0(\Gamma_{z,y}))=e^{ieA_0\cdot(z-y)}$ for every contour inside $\square$, $G_i(\square,\partial\lambda_0)=U(\lambda_0)G_i(\square,0)U(\lambda_0)^*$ and $Q_i(\partial\lambda_0)=U(\lambda_0)_{\mathrm{coarse}}Q_i(0)U(\lambda_0)^*$: \cite[(2.69)--(2.74)]{Ba2}; legitimate because $\square$, of diameter $\le8r_i+s_b$, does not wrap around (on the cube, no condition). Write $\phi'\defeq U(\lambda_0)^{-1}\phi$. (5) \emph{Reproduction of constants:} by \cite[(2.75)]{Ba2}, $a_iG_i(\square,0)Q_i^*1=1-m^2\ell_i^2G_i(\square,0)Q_i^*1$, hence $a_i(G_i(\square,0)Q_i^*\square_1\phi')(x)-\phi'(y)=a_i(G_i(\square,0)Q_i^*[\square_1(\phi'-\phi'(y))])(x)-\phi'(y)a_i(G_i(\square,0)Q_i^*1_{\square\setminus\square_1})(x)-\phi'(y)m^2\ell_i^2(G_i(\square,0)Q_i^*1)(x)$. The first term is $O(p)$: on the convex cube $Q(y,s_i/2)\subset\Lambda_{-1}$ the restriction $|D\phi|\le c_1p$ gives $|\phi'(z)-\phi'(y)|\le O(p)|z-y|$ and the kernel decays exponentially. The second: $\square\setminus\square_1\subset(Q(y,s_i)\setminus Q(y,s_i/2))\cup(Q(y,s_i/2)\setminus\Lambda_6)$ lies at distance $\ge\min\{s_i/2,\rho_i\}-c_d$ from $x$, so it is $\le c_{\mathrm{Ba}}e^{-\delta_0(\min\{s_i/2,\rho_i\}-c_d)}\Upsilon_i$ by \cite[(2.25)]{Ba1} on the parallelepiped. The third is the mass-defect term of the printed proof, $O(p_i)$ under the mass condition. (6) \emph{Derivatives:} a bond $b\subset B^i(y)$ lies at distance $\ge s_i-c_d\ge R_{\mathrm{Ba}}$ inside $\square$ and $\ge5\rho_i-c_d$ inside $\Omega$, so the printed derivative bounds apply at $b$ and (1)--(5) give \cite[(2.66)]{Ba2} in the form \cite[(2.77)]{Ba2}; no re-entrant edge is approached. \emph{Errors:} $s_i/2\ge\rho_i$ iff $\rho_i\le2r_i$ and $3\rho_i\ge2s_b$, the first by $r(\ell_{i-1})/r_i\le2L$, the second by $R\ge10s_b+10$; so every error is $\le Ca_i\Upsilon_i(1+r_i)^{d-1}e^{-\delta_0(\rho_i-c_d)}$, which is $\le\ell_i^\kappa$ by the last hypothesis; the rest of the printed proof, \cite[(2.69)--(2.77)]{Ba2}, runs on $\square$. Nothing depends on $L$ except through the named constants.
\end{proof}
At $L=5$ the hypotheses read $r(\eps_0^{\mathrm{Ba}})\ge R_{\mathrm{Ba}}+(3+\sqrt3)s_b$, $r(\ell_{i-1})/r_i\le6.81\le10$, $R\ge10s_b+10$: lower bounds on $R$ and caps on $\eps_0^{\mathrm{Ba}}$, none restricting $L$. On the cube the same device (a box of large blocks inside the propagator domain about the collar cubes, with the cap inherited from the previous step) is the King box of D8.

% a01_H0rel part 6x3: the replaced step (d) and the parameter conditions (input from a01_H0rel_p6.tex after a01_p6x2)
\subparagraph{(d) The partition \cite[(3.43)--(3.45)]{Ba2}.} The printed ``regular partition into a lattice of cubes, each a sum of large blocks, with side in $(r,2r)$'' needs, per coordinate, a divisor $c$ of the number $N$ of large blocks with $r/s_b<c<2r/s_b$; at $L=2$ a power of $2$ always works, while at $L=5$ the divisors of $N_i=2MJ/(\ell_is_b)$, a number of the form $2\cdot5^{b}J$, have gaps of ratio up to $\frac52$, so the window can contain no divisor, and with integer sides and even $N$ the printed partition fails exactly when $r_i/s_b$ is an integer. The partition is used in \cite[Sect.~3C]{Ba2} only through four properties: the covering \cite[(3.43)]{Ba2}; the neighbour count $3^d$ in \cite[(3.44)--(3.45)]{Ba2}; the volume bounds $(2r)^d$ in \cite[(3.48)--(3.49)]{Ba2}; and the corridor step of \cite[p.~593]{Ba2}.

\begin{lemma}[the product partition and the cube count]\label{lem:a01:PP}
Let $q_i\defeq\lceil(r_i+s_b)/s_b\rceil$ and let every coordinate carry $N\ge q_i$ large blocks (on King's cube, $N=2MJ/(\ell_is_b)$ large blocks per side in scale-$i$ units, which exceeds $q_i$ for $J\ge J_0$; on a torus, $N$ per period). Cut each coordinate interval into consecutive arcs of between $q_i$ and $2q_i-1$ large blocks (write $N=nq_i+\varrho'$, $0\le\varrho'<q_i$, and give the $\varrho'$ extra blocks to the $n$ arcs, at most $q_i-1$ each), and let $\mathcal Q_i$ be the product partition into boxes. Call $Q\sim Q'$ if their arc indices differ by at most one in every coordinate. Then: \textup{(Q1)} every box is a union of large blocks; \textup{(Q2)} every box is a product of intervals of at most $2r_i+3s_b\le(2+3s_b/R)r_i$ points; \textup{(Q3a)} every box has at most $3^d$ boxes $Q'\sim Q$; \textup{(Q3b)} if $\dist(z,x)\le r_i+s_b$ then $Q(z)\sim Q(x)$. Consequently, with $I_i\defeq\{z:\dist(z,(\Lambda_7^{(i-1)\prime})^c)\ge s_b\}$ and $\mathcal C_i\defeq\{Q\in\mathcal Q_i:Q\cap\Lambda_0^{(i)c}\cap I_i\ne\emptyset\}$, every $Q\in\mathcal C_i$ satisfies $Q\sim Q(x_e)$ for the charge point $x_e$ of some element $e$ of the family $F_i$ of large-field events at scale $i$, hence $|\mathcal C_i|\le3^d|F_i|$; the volume bound $(2r)^d$ becomes $(2r_i+3s_b)^d$ and the corridor bound $22r$ becomes $22r_i+3s_b$, both $O(1)r_i^d$.
\end{lemma}
\begin{proof}
(Q1) Arcs are unions of large-block intervals. (Q2) An arc has at most $(2q_i-1)s_b<2(r_i+s_b)+s_b$ points, and $s_b\le(s_b/R)r_i$. (Q3a) The boxes $\sim Q$ are the $3^d$ neighbours of $Q$ in the product structure. (Q3b) Every arc has $\ge q_is_b\ge r_i+s_b$ points, so if $|z-x|_\infty\le r_i+s_b$ the two arcs in each coordinate are equal or consecutive (otherwise a whole arc lies strictly between and the coordinate distance is $\ge q_is_b+1$). For the count: let $Q\in\mathcal C_i$ and $z\in Q\cap\Lambda_0^{(i)c}\cap I_i$; its large block $\beta$ lies in $Q$ by (Q1) and, having sup-diameter $<s_b$, in $\Lambda_7^{(i-1)\prime}$; since $z\notin\Lambda_0^{(i)}$, $\beta$ is not one of the blocks of \cite[p.~570]{Ba2}, so $\dist(\beta,S_i)\le r_i$ where $S_i$ is the point set of $F_i$; take $z'\in\beta$, $y\in S_i$ with $\dist(z',y)\le r_i$, $y$ a point of an element $e$ with $\dist(y,x_e)\le L-1<s_b$; then $\dist(z',x_e)\le r_i+s_b$ and $Q=Q(z')\sim Q(x_e)$ by (Q3b); count with (Q3a). The printed \cite[(3.45)]{Ba2} and the first inequality of \cite[(3.46)]{Ba2} do not hold as printed (a large block straddling $\partial\Lambda_7^{(i-1)\prime}$ lies in $\Lambda_0^{(i)c}$ with $F_i$ possibly empty); the count above, on the interior $I_i$, becomes the count of D14 over the paid elements through \cref{lem:a01:PP5}; in D14 the boxes of $\mathcal C_{i'}$ for $i'<i$ are carried along by the corridor construction of \cite[(3.48)--(3.50)]{Ba2} with the thickness $7(r_{i-1}+s_b)+Ls_b$ inside Ba{\l}aban's $(9r,10r)$ once $(7+L)s_b<2r_{i-1}$.
\end{proof}
% a01_H0rel part 6x3, continued: the count over the paid elements (input from a01_p6x3.tex after lem:a01:PP)
\begin{lemma}[the count over the paid elements]\label{lem:a01:PP5}
Let the boxes $\mathcal Q_i$, the relation $\sim$, the arcs of $q_i$ to $2q_i-1$ large blocks, the set $\mathcal C_i$ and the family $F_i$ be those of \cref{lem:a01:PP}, with the charge point $x_e$ of an element $e$ as in its proof (every point of $e$ lies within $\dist\le L-1$ of $x_e$). Let $F_i^{\mathrm{paid}}\subset F_i$, and suppose that every $e\in F_i\setminus F_i^{\mathrm{paid}}$ has a \emph{witness} $w(e)\in F_i^{\mathrm{paid}}$ with $\dist(x_e,x_{w(e)})\le r_i+s_b$. Then
\[
|\mathcal C_i|\ \le\ 5^d\,|F_i^{\mathrm{paid}}|.
\]
In D14 the unpaid elements are the near $R_s$ elements, the blocks $y\in R_s^{(i)}$ within $\ell^\infty$-distance $R_0$ of the face layer, and $F_i^{\mathrm{paid}}$ consists of Ba{\l}aban's six families, $R_s$ restricted to its far elements, and the face family; the witness of D7(b) (\cref{lem:a01:F541}(e)) has a point within distance $dR_0+1$ of a point of $e$, so $\dist(x_e,x_{w(e)})\le dR_0+1+2(L-1)\le s_b+L-2<r_i+s_b$ by $s_b\ge dR_0+L+1$ (D7(c)) and $r_i\ge R>L$, and the hypothesis holds.
\end{lemma}
\begin{proof}
Let $Q\in\mathcal C_i$. The proof of \cref{lem:a01:PP} gives $z'\in Q$ and $e\in F_i$ with $\dist(z',x_e)\le r_i+s_b$. Put $e'\defeq e$ if $e\in F_i^{\mathrm{paid}}$ and $e'\defeq w(e)$ otherwise; in both cases $e'\in F_i^{\mathrm{paid}}$ and $\dist(z',x_{e'})\le2(r_i+s_b)\le2q_is_b$, since $q_is_b\ge r_i+s_b$. In each coordinate the arc indices of $Q=Q(z')$ and of $Q(x_{e'})$ differ by at most $2$: if they differed by $3$ or more, two whole arcs, of at least $q_is_b$ points each, would lie strictly between the two coordinates, whose distance would then be $\ge2q_is_b+1$. A box has at most $5^d$ boxes whose arc indices differ from its own by at most $2$ in every coordinate (fewer on a torus with fewer than five arcs in some coordinate). So $Q\mapsto e'$ maps $\mathcal C_i$ into $F_i^{\mathrm{paid}}$ with fibres of at most $5^d$ boxes, which is the claim. For the application: a near $R_s$ element $e$ at the block $y$ has $\dist(y,x_e)\le L-1$; \cref{lem:a01:F541}(e) gives, on a block path of length $\le dR_0$ from $y$ to the face layer, a bond of the gradient family or a block of the face family, that is an element $w(e)$ of a paid family one of whose points lies within distance $dR_0+1$ of $y$ and within $L-1$ of $x_{w(e)}$; the triangle inequality gives $\dist(x_e,x_{w(e)})\le dR_0+2L-1$, and $dR_0+2L-1\le s_b+L-2$ is $s_b\ge dR_0+L+1$.
\end{proof}


The replaced objects of (a)--(d), $r^\flat$, $D^\flat$, $\zeta^{(i)}_\theta$, the box $\square$ of \cref{lem:a01:lem24} and the partition $\mathcal Q_i$, are built from $(L,M,R,\ell_i,\bar n)$ and the set of event blocks alone, so they are the same on the cube as on the torus (\cref{thm:a06:BALA}); they are auxiliary objects of the proofs, not of the statements, and under (b) the all-small-field action changes by $O(\ell_i^\kappa)|X^{(i)}|$, which the statements absorb. The parameter conditions they add are lower bounds on $R$ ($R\ge R^\flat$; $R\ge10(M+1)$, $R\ge M_\theta+2$; $R\ge10s_b+10$; $R\ge6L\sqrt d$ for the covering of D14) and caps on $\eps_0^{\mathrm{Ba}}$, all fixed in Ba{\l}aban's order of choices ($d$, $a$, $\bar n$, $\alpha_0$, $L$, $M$, $R$, $\varrho$, $p$, $b_0$, then $\eps_0^{\mathrm{Ba}}$), so the order is acyclic.


\begin{theorem}[the perturbative bounds on sub-domains]\label{thm:a01:B2}
Fix a scale $k\le K_f$, a union $\Omega$ of $k$-blocks of the cube (any shape, meeting $\DD$ or not), $A$ regular on $\Omega$ with $\theta_k\le\theta_0$, and a renormalised family $\mathcal R^\Omega$ of order $\le13$ of Ba{\l}aban's expansion with the lines $G_{(i)}(\Omega,A)$, $G_k(\Omega,A)$, the legs $\phi^{(k)}_\Omega$, photon legs of the class of \cite[p.~412]{Ba3}, possibly source vertices, possibly one difference line $\delta G_k(\Omega,\Omega_2;A)$ of two Neumann domains as an external pair, with at least one external field or source vertex, and all vertex cubes at depth $\ge R_*\defeq48$ in every propagator domain. Then for one held vertex $x_0$ (or none),
\[
\sum_{\text{orderings, indices, positions}}|E^{(k)}(G)|\ \le\ C\ell_k^{\beta_0}p(\ell_k)^{n_1}(1+K-k)^{p_0}e^{-\delta\dist(V_*)}\Lambda(x_0)\,\Pi_b,
\]
with $\Lambda(x_0)$ as in \cref{prop:a01:TD}(a), $\Pi_b=1$ without a difference line and $\Pi_b=\exp[-\delta_c(\dist(\mathrm{cube}(v),\Xl\setminus\Omega_2)+\dist(\mathrm{cube}(v'),\Xl\setminus\Omega_2))]$ for the ends $v,v'$ of the difference line; the constants do not depend on $k$, $K$, $\vth$, $J$, $\Omega$, $A$ or $x_0$. The same holds for a finite set of such difference lines and difference legs, each contributing its own factor, with the constant multiplied by $C_1^{|S|}$.
\end{theorem}
\begin{proof}
The proof of \cref{prop:a01:TD}(a) is run on $\mathcal R^\Omega$ with the following changes. (1) \emph{Localisation.} Every line and leg of the family is localised on a Dirichlet box $Q_x$ of large blocks around its vertex cube, inside $\Omega$ at depth $\ge R_*$: the difference between the operator on $\Omega$ and the operator on $Q_x$ is a cut-off and commutator argument of the type of \cref{lem:a01:FLD} (a smooth cut-off at the Neumann faces of $Q_x$, the commutator bounded by \cref{lem:a01:A5}, the result bounded by \cref{lem:a01:A4}(c)), exponentially small in the depth, and the box-vs-whole-cube difference near $\DD$ is \cref{lem:a01:LOC}; so every line at a vertex of depth $\ge R_*$ is the line of the box up to $Ce^{-\delta R_*}$ times the box line, with the Combes--Thomas rate, free of Ba{\l}aban's $M$. (2) \emph{Zero-background identities on the box.} The five identities used in the organisation of \cite[p.~330]{K2} (the free counterterm, the bulk rewriting of the renormalised families, the constant-datum identity \cref{lem:a01:RC}, the slice identity, and the Ward-type cancellation of the zeroth photon moment, which on a box is exact since a constant gauge field is pure gauge there) hold on the Dirichlet box by the odd-image representation of Step~1 of \cref{lem:a01:FL}, and at a regular background by gauging away the constant part of $A$ on the box (gauge covariance, \cref{app:a01:weights}) and expanding to first order in the remainder, whose coefficients are $O(\theta_k(1+|x-x_0|))$ (\cref{lem:a01:FL}, Step~3); the first-order photon tadpole at a background is bounded by the same expansion. (3) \emph{Remainder vertices.} Each boundary-remainder vertex is a layered vertex of the class of \cref{prop:a01:TD}(a) conjugated by the gauge transformation of (2), plus an ordinary vertex of degree raised by $D'+2\ge1$, plus an any-degree term with the small factor of (1); conjugation does not change absolute values, so the degree bookkeeping of \cite[Sect.~2]{Ba3} is unchanged and the extra vertices only raise the containing stages. (4) \emph{Lines near the face, difference lines.} The decay along paths of the lines on $\Omega$ is \cref{lem:a01:A23,lem:a01:A4}; the H\"older clause of a difference line near $\DD$ comes from the proof of \cref{lem:a01:A5}, cases (3)--(4), applied to $w=\Mop(\Omega_2)^{-1}g$ with $g$ on the cut shell. The difference line $\delta G_k(\Omega,\Omega_2)=G_k(\Omega)-G_k(\Omega_2)$ is $G_k(\Omega)Y G_k(\Omega_2)$ with $Y$ the bond terms of the cut, so its kernel between two points carries the factor $e^{-\delta_c(\dist(\cdot,\Xl\setminus\Omega_2)+\dist(\cdot,\Xl\setminus\Omega_2))}$ by Combes--Thomas on both sides, which is $\Pi_b$. (5) \emph{Telescoping.} A family some of whose lines are written $o=o'-\Delta o$ is the finite sum $\sum_S\mathcal R^{(S)}$ over the sets $S$ of decomposed objects, $\mathcal R^{(S)}$ carrying $\prod_S(-\Delta o)$ and $\prod_{O\setminus S}o'$; $\mathcal R^{(\emptyset)}$ is the family with every object replaced by its counterpart on $\Omega$, with the same counterterms, and for $S\ne\emptyset$, $\mathcal R^{(S)}$ is the renormalised family of the cut graphs $G\setminus S$, with every subtraction of the divergent subgraphs of $G\setminus S$ and the same counterterm values, the divergent subgraphs containing a line of $S$ carrying no subtraction: this is the regrouping \cref{hyp:k01:RENorb}, and no subtraction is lost. Each $\mathcal R^{(S)}$ is then a family with $|S|$ admissible external pairs, bounded by (1)--(4) with the factor $\Pi_S=\prod_{\sigma\in S}C_\sigma e^{-\delta_\sigma(\dots)}$; the constant per pair is $C_1=C_1(d,L,\bar n)$. (6) \emph{Depth.} At every consumer (D5, D6, D7, D10 below) the vertex cubes lie in a $\Lambda_7$-level set and the propagator domains are $\Lambda_2$-level, five nested layers apart (\cite[p.~411]{Ba3}, \cite[(2.53), (2.56), Prop.~2.6]{Ba2}), so their depth in every propagator domain is $\ge\lfloor R\rfloor-1\ge51>R_*$ once $R\ge\max(52,\lceil R_{\mathrm{Ba}}\rceil+1)$, where $R_{\mathrm{Ba}}$ is Ba{\l}aban's depth constant of \cite[(2.19)]{Ba4}.
\end{proof}

\begin{proposition}[changes of propagator domain]\label{prop:a01:dGD}
At every scale $k\le K_f$ of the upper bound, the domain-change statements of \cite[(2.18)--(2.50), (2.108)--(2.115)]{Ba2} hold on the cube with the weighted flow: (i) the comparison of a Dirichlet-conditioned covariance $C_\Lambda$ with the covariance of the Neumann domain $B^k(\Lambda_2^{(k)})$, $C_\Lambda-C|_{\Lambda\Lambda}=C_\Lambda X_{\Lambda\Lambda^c}C_{\Lambda^c\Lambda}$, is $O(M_{k+1}e^{-cr(\ell_k)})$ in the external-field norm; (ii) the scalar translation \cite[(2.23)]{Ba2} and its replacement \cite[(2.25)]{Ba2} differ by $O(M_1e^{-\delta r(\eps)})=O(\eps^{\kappa'})$ for every $\kappa'$; (iii) the change \cite[(2.108)]{Ba2} of the block operators $\Delta^{(k)}$ between $B^k(\Lambda_2^{(k-1)\prime})$, $B^k(\Lambda_2^{(k-1)\prime}\cap\Lambda_5^{(k)c})$ and $B^k(\Lambda_2^{(k)})$ obeys \cite[(2.109)]{Ba2}, $|h_k(x,x')|\le Ce^{-c\min(r(\ell_k),4r(\ell_{k-1})/L)}e^{-\delta|x-x'|}$; (iv) the Neumann cuts along $\partial\Lambda_6$, $\partial\Lambda_7$ inside the Dirichlet-restricted covariances \cite[(2.38)--(2.40)]{Ba2} cost $O(\eps^{\kappa'})|\partial\Lambda_5|$, resp.\ $O(\eps^{\kappa'})|\Lambda_5|$; (v) the exact recursion \cite[(2.41)]{Ba1} on the Neumann domain $B^k(\Lambda_2^{(k)})$ holds (it uses only the uniform coarse step); (vi) the restricted fluctuation covariance $C^{(k)}_{\Lambda_5^{(k)}}$ of \cite[(2.115)]{Ba2} and the Neumann hole of \cite[(2.46)--(2.47)]{Ba2} cost $O(\ell_{k-1}^{\kappa'})|\Lambda_6^{(k-1)}|$; (vii) the photon statements are those of the torus by \cref{lem:a01:PE}(c). The constants depend on $(d,L,a,\Wt,c_V)$ and King's and Ba{\l}aban's constants only, not on $J$, $K$, $k$, $\eps$, the domains (any unions of blocks in the classes where the objects are defined, the face met or not, re-entrant edges allowed) or the field.
\end{proposition}
\begin{proof}
Each statement is an identity of resolvents (the bond cut $C_\Lambda-C'_\Lambda=-C_\Lambda Y_{\Lambda\Lambda}C'_\Lambda$, the restriction identity above, the Schur complement for the hole) followed by the Combes--Thomas decay of \cref{lem:a01:A23}(a) on each factor and the coercivity \cref{lem:a01:A1} on every domain (which is what replaces Ba{\l}aban's $\gamma_0'$; the step-$0$ lower bound for every field is the Neumann gap with the weight, as in the proof of \cref{lem:a01:A1}), the sup bounds of \cref{lem:a01:A4} and the legs' norms of \cref{lem:a01:legs} on the domains involved; the changes relative to the torus are constants only (the factor $\Wt$ from (Q5), (Q6) wherever a block average is bounded, and $\gamma_0''$ of \cref{lem:a01:A23}(c) in place of Ba{\l}aban's $\gamma_0'$). The distances $r(\ell_k)$, $4r(\ell_{k-1})/L$ between the sets compared are Ba{\l}aban's region geometry, unchanged. The photon part is \cref{lem:a01:PE}(c).
\end{proof}

